%% Original-source obligations and explicit repairs; not binary verification.
\begin{annexitem}{annex:p12-grasp-causes-sound}{thm:p12-grasp-causes-sound}{Termination and semantic sufficiency of GRASP causes}
\begin{proof}
Order the finite graph by the time at which assignments were recorded.
Every selected unit reason uses earlier assignments, and a conflict uses
already present assignments.  Recursion through same-level antecedents
therefore strictly decreases this order.  Lower-level antecedents are
retained as leaves, so the recursion terminates even for a same-frame FDA.

Induct on this order.  A vertex returned as its own assignment has that
value in every valuation satisfying the cause.  Otherwise the causes contain
the lower-level antecedents and all causes of same-level antecedents.
By induction a model satisfying them makes every selected reason antecedent
true.  The reason is an input consequence and forces the recorded value.
At the conflict vertex this argument makes every literal of its conflict
clause false, contradicting that input consequence.  Hence no input model
extends the conflict causes; their negation is an input implicate.  An empty
cause set excludes every model and gives the empty implicate.

This reasoning does not assume FDA reasons sound without justification.
Induct also on learned-clause creation.  The first diagnosis uses original
reasons and yields an input implicate by the graph argument.  Its FDA is
then a unit consequence of that proved implicate.  Each later diagnosis
uses originals and earlier proved implicates, so its new clause is again
an input implicate.  This covers any finite sequence of FDA explanations;
they are expanded as implications rather than treated as fresh decisions.
Fixed external query assumptions may remain as cause leaves: the resulting
clause is conditional on them and does not assert that the assumptions are
themselves consequences of the input.
\end{proof}
\end{annexitem}

\begin{annexitem}{annex:p12-grasp-standard-search}{thm:p12-grasp-standard-search}{The standard GRASP search with a declared root convention}
\begin{proof}
Recorded clauses preserve all input models by the cause lemma.  A unit step
preserves every model extending the current prefix: its sole free literal
must be true.  A root conflict excludes all models extending \(A_0\).
Fixed root facts and assumptions are never erased.  A nonconflicting total
return satisfies every original clause, since a false clause would have
been reported by completed Deduce.  If a partial prefix already satisfies
the database, arbitrary completion also gives a model.

At a positive frame \(d\), the new decision is its only elective assignment.
Standard causes expand current-level implications, stopping at that decision
or earlier assignments.  If the decision remains in the causes, their largest
level is \(d\).  Erasing this attempt frees its literal and leaves all
earlier cause literals false in the recorded clause.  That clause is unit
and forces the opposite value: the FDA.  Every model of the earlier prefix
must obey it; this is not a second freely chosen decision.

The original elective decision has now been removed.  Assignments of the
second current-frame attempt are FDAs or their consequences, whose backward
explanations end in earlier frames.  A new local conflict therefore has
cause maximum below \(d\).  If a deeper call returns a conflict targeted
at \(d\), the present diagnosis expands its current-frame FDA causes in
the same manner and again obtains a smaller target.  A frame has at most
its initial decision attempt and its forced opposite attempt.  This progress
fact uses standard diagnosis and retained reasons; arbitrary substitute
engines do not acquire it merely by using the name FDA.

When the maximum cause level is \(\beta<d\), the whole cause assignment
is already in the prefix through \(\beta\).  Every extension falsifies
its recorded implicate.  Passing the conflict over deeper frames therefore
discards only prefixes having no model.  At the target frame diagnosis
produces its FDA or passes a smaller target upward.  Strict decreases among
nonnegative levels cannot continue indefinitely without reaching such a
frame or the checked root.

For termination, induct on the number of unassigned original variables at
frame entry.  Each deduction phase is finite; each decision assigns a new
variable.  Every deeper call consequently has fewer unassigned variables.
There are at most two attempts at a frame, so induction gives finite return
of its children and of the frame itself.  With no variables left the checked
terminal state returns a model or an already detected conflict.  Learned
clauses introduce no variables and are input implicates, so database growth
does not invalidate this decreasing measure.

For completeness suppose a model extends the entry prefix.  If it has the
initial decision value, sound propagation preserves it and the smaller-
variable induction gives success in that branch.  If the first attempt
fails while some entry-prefix model still exists, its target cannot be below
this frame, since such a target would exclude every entry-prefix extension.
Its local FDA forces the opposite value and preserves all remaining models.
The second attempt then succeeds by the same induction.  Thus failure
excludes every entry-prefix model.  At the root these are precisely the
models extending \(A_0\).  This proves both return specifications and
termination without importing restarts, selective deletion or an unexplained
zero-based wrapper erasure convention.
\end{proof}
\end{annexitem}

\begin{annexitem}{annex:p12-grasp-uip-links}{thm:p12-grasp-uip-links}{Soundness of the corrected native UIP links}
\begin{proof}
In forward UIP order, an earlier UIP dominates a later one: a decision-to-
later-UIP path bypassing the earlier one could be extended to the conflict,
contradicting domination.  Acyclic assignment order prevents that extension
from returning through the earlier node.

For adjacent \(u_i,u_{i+1}\), expand reasons backward from the downstream
\(u_{i+1}\), stopping at \(u_i\) and retaining older-level assignments.
No other current decision leaf can be reached while bypassing \(u_i\),
by domination.  Finite graph order proves termination.  Induction on the
expanded reasons shows that every input model satisfying the returned
antecedent assignments gives \(u_{i+1}\) its recorded value.  Their
negation, disjoined with the literal asserting that downstream value, is
therefore an input implicate.  Expanding from \(\kappa\) with stop \(u_k\)
similarly makes conflict inevitable under its returned assignments; negating
them gives the terminal conflict implicate.

The stopping assignment remains in the antecedent part with its false
literal.  The appended true literal must concern the downstream output,
not the stopping assignment.  Appending that stopping assignment's own
true literal would make a tautology.  Reversing the cause arguments would
attempt to explain an upstream vertex by a downstream one and would not
prove the required implication.  This establishes the corrected endpoint
order rather than merely renaming the printed Eq.~(9).

In the native example, the support of \(x_4=1\) relative to \(x_1=1\)
is \(\{x_1=1,x_9=0\}\), giving \(\neg x_1+x_9+x_4\).
The conflict support relative to \(x_4=1\) is
\(\{x_4=1,x_{10}=0,x_{11}=0\}\), giving
\(\neg x_4+x_{10}+x_{11}\).  These are obtained from the actual
\(\omega_1\)--\(\omega_6\) reasons and justify the stated two-clause
example without treating the faulty index formula as a proof.
\end{proof}
\end{annexitem}

\begin{annexitem}{annex:p12-grasp-complete-product}{thm:p12-grasp-complete-product}{Conflict-free decisions with all prime implicates}
\begin{proof}
For a consistent partial assignment \(B\), its blocking clause consists of
the literals opposite to its assigned values.  If \(B\) has no model
extension, that clause is an implicate: a model falsifying it would extend
\(B\).  Choose an inclusion-minimal entailed subset among its finitely many
subsets.  It is a prime implicate, and all its literals are false under
\(B\).  Hence a nonconflicting prefix on \(P(\varphi)\) has a model
extension.

Let such a prefix \(A\) also be unit-saturated, and choose a value \(v\)
for an unassigned \(x\).  If \(A\cup\{x=v\}\) had no model extension,
the preceding argument would give a prime implicate contained in its
blocking clause.  It cannot be contained in the blocking clause of \(A\),
or it would already be false there.  It therefore contains the literal
opposite to \(x=v\), with all others already false under \(A\).
Before the decision it was unit and would have forced the opposite value,
contrary to \(x\) still being unassigned at saturation.  Thus either
decision value retains a model; sound propagation preserves it and cannot
conflict.  Repeat at every decision.

For unsatisfiable \(\varphi\), the empty clause is entailed and prime.
Every nonempty implicate has it as a proper entailed subset and is not prime,
so the complete product is exactly the empty clause.  For the identically
true function, every non-tautological clause has a falsifying valuation and
is not entailed; the prime product is empty.  These cases include an empty
variable set.  The needed direction is that a prime subclause
\(\omega\subseteq\psi\) entails \(\psi\), not conversely.
No small-size bound is used, and externally contradictory initial assumptions
are not a compatible prefix of the satisfiable case.
\end{proof}
\end{annexitem}

\begin{annexitem}{annex:p12-minisat-original-reason-bridge}{thm:p12-minisat-original-reason-bridge}{The nonclausal reason-to-clause bridge}
\begin{proof}
Let \(F\) conjoin the permanent constraints.  Each constraint \(K\) is
a conjunct, so \(F\models K\).  Its propagation-reason contract gives
\(F\models(\bigwedge R\to p)\), exactly the clause made from the
complements of \(R\) and \(p\).  Its conflict contract similarly entails
the clause of those complements alone.  This works for any Boolean
predicate constraint, not just a syntactic clause.  The contract is an
explicit premise, not a property conferred by a class name.

The support literals are true in the relevant saved state and precede their
propagation.  Correct undo restores that state, making the explanation graph
well founded.  Ordinary pivot resolution preserves entailment, as proved in
Annex~\ref{annex:p12-cdcl-uip-local-sound}.  Induction over callback expansion
and learned-clause creation proves that every learned clause is a permanent-
constraint consequence.  The general first-UIP graph proof applies only
with its actual chronological/current-level/propagation premises; the original
constraint need not itself be a clause.  Omitting level-zero support is sound
because the permanent base entails those root facts: resolution with their
unit consequences removes their false complements.  Temporary assumptions
are not such permanent facts and may not be omitted merely because their
level is called a query root.

Event coverage guarantees that every violated permanent constraint at a
total assignment would have exposed a conflict.  A total conflict-free return
therefore satisfies all of them.  Without coverage, a sound callback never
invoked on its necessary event could not justify SAT.  At a query-root
conflict, its supporting assignments follow from permanent constraints and
fixed query assumptions.  The conflict clause is an input consequence, but
all its literals are false under those query consequences, excluding every
input model obeying the assumptions.  This is a relative refutation, not
verification of missing callbacks, watch registrations, undo or ignored
stronger-simplifier results.
\end{proof}
\end{annexitem}

\begin{annexitem}{annex:p12-minisat-original-query}{thm:p12-minisat-original-query}{Root simplification and assumption-relative conclusions}
\begin{proof}
Every \(F\)-model satisfies \(G\).  Under \(G\), retained replacements
are equivalent to their original constraints and deleted constraints are
already true.  Thus every \(F\)-model satisfies \(G\land F'\).
Conversely a \(G\land F'\)-model satisfies every old constraint, the
retained ones by equivalence under \(G\), the deleted ones because \(G\)
makes them true.  Hence it satisfies \(F\).  This proves model equality
for the represented conjunction including root facts, not for \(F'\) alone.

Temporary assumptions \(H\) intersect that same model set with valuations
satisfying \(H\).  Sound propagation and the reason bridge preserve all
those models.  A total returned model satisfies \(F\land H\); a sound
query-root conflict excludes all such models.  Cancelling the temporary
assumptions leaves the permanent representation unchanged.  Learned clauses
remain globally reusable only if their derivation kept temporary assumption
literals in its reason clauses rather than treating them as permanent facts.

The example \(F=x+y\), \(H=\neg x\land\neg y\) proves that query
UNSAT need not be base UNSAT.  A budget exit interrupts search and supplies
neither a contradiction nor a total model, so its undefined result licenses
neither conclusion.  The pre-solve simplification check is also needed to
distinguish a permanent conflict from a temporary one in the source object
having no persistent inconsistency flag.  Ignoring a false constructor or
simplification return violates the usable-base premise.
\end{proof}
\end{annexitem}
